\scanpage{135}
\mspage{14}
$(A\bar R)\cap R=A$. Now it is obvious that if $A$ is principal then so is $A\bar R$, so assume conversely that $A\bar R$ is principal and let $\bar y$ be a generator of $A\bar R$. Then it is obvious that if $\bar y=0$ then $A=0$, so assume that $\bar y\ne0$ and let $q$ be the unique integer for which $\bar y\in\bar M^q$ and $\bar y\notin\bar M^{q+1}$ so that $A\bar R\subset\bar M^q$ and $A\bar R\not\subset\bar M^{q+1}$. Since $M^{q+1}\bar R=\bar M^{q+1}$, there exists $y$ in $A$ such that $y\notin M^{q+1}$. Since $\bar M^{q+1}\cap R=M^{q+1}$, $y\notin\bar M^{q+1}$. Now $y=\bar yz$ with $z$ in $\bar R$. Since $\bar y=\bar M^q$ and $y\notin\bar M^{q+1}$, $z$ must be a unit in $\bar R$ and hence $A\bar R=y\bar R$ so that $A=(A\bar R)\cap R=((yR)\bar R)\cap R=yR$.

Finally take $y$ in $K\cap\bar R$, then $y=z/t$ with $z\in R$ and $0\ne t\in R$ so that $z=ty\in(t\bar R)\cap R=tR$ so that $y=z/t\in R$. Therefore $K\cap\bar R\subset R$ and hence $K\cap\bar R=R$ and $K\cap\bar M=(K\cap\bar M)\cap\bar R=\deleted{R=}(K\cap\bar R)\cap\bar M=R\cap\bar M=M$.

LEMMA 3. (A form of Weirstrass Preparation Theorem). Let $(R,M)$ be a local domain of dimension $n>1$ such that $R$ and $R/M$ have the same characteristic so that $R$ has a coefficient field $k$; let $x_1,\ldots,x_n$ be a basis of $M$ (then $R=k[[x_1,\ldots,x_n]]$). Let $y$ be a non-zero element of $R$ and let $b=\lambda_R(y)$. If the line $x_2=\cdots=x_n=0$ is non-tangential to $y$ then there exist unique elements $\deleted{t}\inserted{y}_1,\ldots,\deleted{t}\inserted{y}_q$ in $k[[x_2,\ldots,x_n]]$ and a unique unit $\eta$ in $R$ such that $y=y'\eta$ where $y'=x_1^b+\deleted{y_1^b}y_1x_1^{b-1}+\cdots+y_{\deleted{q}\inserted{b}}$. Now let $z$ be a non-zero element of $k[[x_2,\ldots,x_n]][x_1]$ and let $t$ be a non-unit factor of $z$ in $R$, and let $d=\lambda_R(t)$. If the line $x_2=\cdots=x_n=0$ is tangential to $z$ then it is also non-tangential
\mspage{134}
